\section{Mikroműholdak bemutatása}
Az űrtechnológia fejlődésével egyértelművé vált az ipar számára, hogy olyan szakemberekre van szükség a nagy projektekhez, akik már gyakorlottak legalább minimális szinten a műholdfejlesztés folyamatainak egy részében. A 2000-es évek elején Robert Twiggs, amerikai professzor javaslatára egységes méret-rendszert hoztak létre a kis műholdak számára. Így alakult ki a Twiggs-elv, mely alapján az egységnyi kis műhold mérete $10cm\times10cm\times10cm$, tömege maximum $1kg$. A kialakításból adódóan ezeket CubeSat-nek kezdték el becézni. Mivel egységnyi méretet szabtak meg, így könnyen skálázhatóvá vált, hogy mekkora mikroműholdat lehet készíteni adott célra.

Szakdolgozatom témájaként egy ilyen CubeSat műhold fedélzeti számítógépét kezdtem el fejleszteni hardvertől szoftverig. A folyamat során inspirált az elérni kívánt cél, a már meglévő magyar sikerek és a mikrovezérlők programozása iránti kíváncsiságom is~\cite{masat-origo}.

\subsection{Modulok egy CubeSat műholdon}
Egy CubeSat műholdnak több felépítő egysége van, melyek közül kritikus elemek a kommunikációs modul, az energiaellátásért felelős napelemek és tápegység, illetve a mindent összekötő fedélzeti számítógép (OBC = On Board Computer).

Mindezek mellett szükséges minimum egy, valamiféle kísérletet vagy egyéb funkcionalitást megvalósító egység is annak érdekében, hogy ne csak kommunikálni lehessen a műholddal, hanem 
legyen valami haszna is. %Úgy írnám, hogy rentábilis legyen. Szeretem a kőművesnyelvet és művelem is, egy szakdoga ennél egy hajszállal (nem többel) fennköltebb.
Mivel szakdolgozatomban az OBC-vel foglalkoztam, így azon modult fejtem ki jobban a következőekben.

%Ide jobban ki lehetne fejteni, akár egy szép ábrával, hogy mik vannak egy műholdban, mikkel kommunikál egy OBC.

\Aref{fig:obcfeladatok}. ábra foglalja össze a legfőbb feladatokat.

\begin{figure}
        \centering
        \begin{adjustbox}{width=\textwidth} % Scale the forest to fit text width
\begin{forest}
for tree={
  grow'=south, % Children grow downwards
  align=center, % Center-align text in nodes
  edge path={
    \noexpand\path [draw, \forestoption{edge}] (!u.parent anchor) -- (.child anchor)\forestoption{edge label};
  },
  inner sep=2pt, % Padding inside nodes
  l sep+=0.5cm, % Vertical spacing between levels
  s sep+=0.5cm % Horizontal spacing between siblings
}
[On-Board Computer
  [Kommunikáció
    [Csomagküldtés
      [Csomagküldtés
        [Csomagküldtés]
      ]
      [Csomagküldtés
        [Csomagküldtés]
      ]
    ]
    [Trials during the test
      [Appearance of the stimuli
        [Time series for each epoch]
      ]
      [Answer to the stimuli
        [Time series for each epoch]
      ]
    ]
  ]
  [Recordings from the  test
        [\{\ldots\}]
  ]
]
\end{forest}
\end{adjustbox}
        \caption{Az OBC feladatai}
        \label{fig:obcfeladatok}
\end{figure}



\section{Fedélzeti számítógép}
A fedélzeti számítógép a lelke minden űreszköznek, akárcsak az újabb gépjárműveknek~is.

Az OBC időzíti a perifériák működését, ellenőrzi azok működőképességét és felel a műhold moduljainak összhangban tartásáért. Mindezek mellett adatokat gyűjt, melyek alapján optimalizálni lehet a működést, későbbi küldetésekhez szerezhetnek a fejlesztők/kutatók további információkat.

\subsection{Fejlesztéshez szükséges eszközök}
A fejlesztés két nagyobb fázisra bontható szét. Az első fázisban a hardvert kellett megtervezni, melyhez a KiCad~8.0~\cite{kicad} ingyenes tervezőprogramot használtam. Miután megérkeztek az alkatrészek és a lész nyomtatott áramkör, kézi forrasztással helyére illesztettem az alkatrészeket. Az áramkör \aref{fig:kapcsolasi_rajz}.~ábrán látható.

\subsubsection{Nyomtatott áramkör fejlesztése}

